\chapter{Limitations}

Despite the robust capabilities of Project Satya in identifying and analyzing misinformation, several inherent constraints remain. Acknowledging these limitations provides transparency into the system's operational boundaries and highlights areas for future iteration.

\section{Local Hardware Constraints and Inference Latency}
The application depends heavily on the execution of embedded Large Language Models (LLMs) and the continuous embedding of incoming claims to query the Vector Database. As such, the system is fundamentally bottlenecked by the available local compute resources, specifically GPU VRAM and CPU thread availability. 

Unlike API-driven architectures that offload compute to cloud providers, processing heavy multimodal payloads (like deepfake video inference and concurrent text generation) locally introduces significant latency. High-volume usage environments could experience throttling or out-of-memory (OOM) errors, rendering the application temporarily unable to securely analyze claims without expanding the underlying hardware infrastructure.

\section{Dataset and Hallucination Risks}
A critical limitation of all Generative AI models is the potential for hallucination—generating plausible but incorrect reasoning. Although strict prompting guidelines are implemented to demand an "Unverified" verdict when internal context is lacking, there remains a persistent risk of the model confidently associating incorrect evidence with a claim. Therefore, the architectural design explicitly mandates that Satya operates as a \textit{Generative Verification Assistant} rather than an absolute source of truth. Human oversight is still necessary to validate the citations generated by the LLM.

\section{Multimodal Extraction Limitations}
The image verification pipeline currently hinges on standard Optical Character Recognition (OCR) tools to extract visible textual context. In scenarios where a manipulated image contains no text—but instead alters the underlying visual semantics (e.g., swapping a face without changing captions)—the OCR pipeline may not provide enough context for the LLM to successfully flag the manipulation compared to dedicated architectural Vision Transformers.
